\documentclass{report}
\usepackage{amsmath,amssymb}
\setlength{\parindent}{0mm}
\setlength{\parskip}{1em}
\begin{document}
\begin{center}
\rule{6in}{1pt} \
{\large Steffen Muenzenmaier \\
{\bf A Comparison of Finite Element Spaces for $H(div)$-Conforming First-Order System Least Squares}}

1ECOT \\ Engineering Office Tower \\ 1111 Engineering Drive \\ Boulder CO 80309
\\
{\tt steffen.muenzenmaier@colorado.edu}\\
Christopher Leibs\\
Thomas Manteuffel\end{center}

First-order system least squares (FOSLS) is a commonly used technique in
a wide range of physical applications. FOSLS discretizations are
straightforward to implement and offer many advantages over traditional
Galerkin or saddle point formulations. Often these problems are
formulated in $H(div)$ spaces and $H(div)$-conforming elements are used.
These elements have lesser regularity assumptions than the commonly used
$H^1$-conforming elements and are therefore believed to be more suited
for singular problems arising in many applications. This talk will
compare the approximation properties of the $H(div)$-conforming
Raviart-Thomas and Brezzi-Douglas-Marini elements to $H^1$-conforming
piecewise polynomials in a $H(div)$-setting. Furthermore a
$H^1$-formulation for these problems will be used and compared to the
$H(div)$-formulation. For the comparison typical Poisson/Stokes problems
are examined and singular solutions will be addressed by adaptive
refinement strategies.


\end{document}
